% ============================================================
% Springer LLNCS Format - Conference Paper
% AI-Integrated Multi-Level Parking Building Management System
% with Automatic License Plate Recognition (AI-LPR)
% SE1820 Group 3 - FPT University
% ============================================================
\documentclass[runningheads]{llncs}

\usepackage{graphicx}
\usepackage{amsmath}
\usepackage{booktabs}
\usepackage{hyperref}
\usepackage{url}
\usepackage{xcolor}
\usepackage{array}
\usepackage{multirow}
\usepackage{pifont}          % \ding{51} = ✓, \ding{55} = ✗
% Uncomment for blue URLs (Springer eBook style):
% \renewcommand\UrlFont{\color{blue}\rmfamily}

\begin{document}

% ============================================================
% TITLE
% ============================================================
\title{An AI-LPR-Integrated Multi-Level Parking Building Management System with Real-Time Slot Allocation and Automated Billing}

\titlerunning{AI-LPR Multi-Level Parking Management System}

% ============================================================
% AUTHORS
% ============================================================
\author{
  Ho Minh Nghia\inst{1} \and
  Le Quoc Huy\inst{1} \and
  Pham Trung Tin\inst{1} \and
  Dam Manh Dung\inst{1} \and
  Bui Quang Vinh\inst{1}
}

\authorrunning{H. M. Nghia et al.}

\institute{
  FPT University, Ho Chi Minh City, Vietnam \\
  \email{\{nghiahmse182398, huyse182088, tinptse182457,
         dungdmse182094, vinhbqse182589\}@fpt.edu.vn}
}

\maketitle

% ============================================================
% ABSTRACT
% ============================================================
\begin{abstract}
Urban multi-level parking facilities face persistent challenges
in managing high vehicle turnover, preventing reservation
conflicts, and minimizing driver search time across multiple
floors. Existing systems predominantly rely on manual entry,
rule-based allocation, or isolated sensor approaches that fail
to address the compound complexity of multi-floor environments.
This paper presents an integrated Parking Building Management
System (PBMS) that combines \emph{Automatic License Plate
Recognition} (AI-LPR) at entry/exit gates with a real-time
slot-allocation engine and an LLM-assisted smart parking
assistant. The AI-LPR pipeline pairs a fine-tuned
YOLOv8-based vehicle detector with a Convolutional
Recurrent Neural Network (CRNN) OCR module, achieving a
target plate-character accuracy above 85\% under standard
lighting. The slot allocator employs a heuristic optimization
strategy that minimises driver travel distance, reduces average
search time by at least 20\%, and maintains a reservation
success rate above 95\%. A Google Gemini-backed conversational
assistant further guides drivers to optimal slots in natural
language. The system is delivered as a full-stack web and
mobile application (Next.js, React Native, Node.js,
MongoDB) with VNPay automated billing. Experiments across
low-, medium-, and high-load scenarios confirm that the
proposed framework surpasses rule-based and sensor-only
baselines on all key performance indicators while keeping
gate-entrance latency below 30~seconds.

\keywords{Multi-level parking \and Automatic license plate
recognition \and YOLOv8 \and Real-time slot allocation \and
Smart parking assistant \and LLM}
\end{abstract}

% ============================================================
% 1. INTRODUCTION
% ============================================================
\section{Introduction}

Rapid urbanisation has driven vehicle ownership to record
levels, creating an acute shortage of organised parking
capacity in commercial and public buildings. Studies show that
drivers in dense urban areas spend an average of 17 minutes
per trip searching for parking, contributing significantly to
inner-city congestion, fuel waste, and carbon emissions
\cite{ref_shoup2011}. Multi-storey parking buildings
(multi-level facilities) are the most space-efficient response
to this demand, yet their operational complexity outpaces the
management tools available to operators.

Contemporary smart parking literature focuses on one or two
sub-problems in isolation: camera-based occupancy detection
\cite{ref_padmavathi2024}, Automatic License Plate Recognition
(ALPR/LPR) \cite{ref_ieee_multilevel2024}, reservation
workflows \cite{ref_plosone2025}, or predictive availability
\cite{ref_dhopre2024}. Few systems address the full pipeline
from gate identification to floor-aware allocation to
online payment in a single deployable architecture.

This paper closes that gap with four primary contributions:
\begin{enumerate}
  \item An end-to-end \textbf{AI-LPR} module integrating
        YOLOv8 detection with a CRNN-OCR reader, enabling
        touchless check-in and check-out at multi-level gates.
  \item A \textbf{real-time, reservation-aware slot allocator}
        that handles concurrent bookings, conflict resolution,
        and cross-floor load balancing.
  \item An \textbf{LLM-powered parking assistant} (Google
        Gemini with Function Calling) providing
        natural-language slot guidance to mobile users.
  \item A \textbf{comprehensive evaluation framework}
        covering occupancy accuracy, gate latency,
        time-to-park, and user satisfaction across three
        load scenarios.
\end{enumerate}

The remainder of this paper is organised as follows.
Section~\ref{sec:related} reviews related work.
Section~\ref{sec:system} describes the system architecture.
Section~\ref{sec:ai} details the AI components.
Section~\ref{sec:eval} presents the experimental setup and
results. Section~\ref{sec:discuss} discusses findings, and
Section~\ref{sec:conclusion} concludes.

% ============================================================
% 2. RELATED WORK
% ============================================================
\section{Related Work}
\label{sec:related}

\subsection{Camera-Based Occupancy Detection}

Padmavathi and Selvaraj~\cite{ref_padmavathi2024} propose an
Ensemble Deep Network (EDN) with Multi-scale Retinex
preprocessing and an Advanced Pelican Optimization Algorithm
(APOA) for parking occupancy detection, achieving 96.6\%
accuracy on the CNRPark dataset. While highly accurate in
single-floor settings, the method does not address
multi-floor coordination or gate-level recognition.

\subsection{Multi-Level Parking Automation}

Chang et al.~\cite{ref_chang2026} present a Pipeline-based
Traffic Analysis System (PTAS) combining DuCRG-YOLOv11n
object detection (91.4\% precision at 68 FPS), ByteTrack
multi-object tracking, and LLM-based semantic analysis via
RAG+Vertex~AI. The system achieves high intersection
throughput but does not target indoor multi-floor reservation
or gate-level plate identification.

The IEEE study on automated multi-level parking
\cite{ref_ieee_multilevel2024} applies a Region-based CNN
for plate recognition at entry gates (98.21\% spatial accuracy)
and couples Model Predictive Control (MPC) with Proximal
Policy Optimisation (PPO) for robotic vehicle transfer to
parking bays. This work validates AI-LPR at the gate but
focuses on robotised rather than driver-navigated facilities.

\subsection{Smart Slot Allocation and Reservation}

The PLOS ONE system~\cite{ref_plosone2025} integrates ML-based
occupancy prediction with Genetic Algorithm or Ant Colony
Optimisation (ACO) to allocate slots at reservation time,
reducing driver search time by 35--40\%. It does not, however,
couple allocation with a live plate-recognition gate or
mobile payment pipeline.

Dhopre et al.~\cite{ref_dhopre2024} propose a web-based
predictive smart parking system using supervised ML on
historical occupancy patterns. Evaluated against software
engineering metrics rather than quantitative ML scores, the
system demonstrates stable UX but lacks a mobile application
and real-time IoT or camera integration.

\subsection{Research Gap}

Table~\ref{tab:gap} summarises the key gaps our system
addresses relative to prior work.

\begin{table}[ht]
\caption{Gap analysis between existing systems and the
proposed PBMS.}\label{tab:gap}
\begin{tabular}{lcccccc}
\toprule
\textbf{Feature} &
\cite{ref_padmavathi2024} &
\cite{ref_chang2026} &
\cite{ref_ieee_multilevel2024} &
\cite{ref_plosone2025} &
\cite{ref_dhopre2024} &
\textbf{Ours} \\
\midrule
AI-LPR gate          & \ding{55} & \ding{55} & \ding{51} & \ding{55} & \ding{55} & \ding{51} \\
Multi-floor aware    & \ding{55} & \ding{55} & \ding{51} & \ding{51} & \ding{55} & \ding{51} \\
Reservation conflict & \ding{55} & \ding{55} & \ding{55} & \ding{51} & \ding{51} & \ding{51} \\
Mobile app           & \ding{55} & \ding{55} & \ding{55} & \ding{55} & \ding{55} & \ding{51} \\
Automated billing    & \ding{55} & \ding{55} & \ding{55} & \ding{55} & \ding{51} & \ding{51} \\
LLM assistant        & \ding{55} & \ding{51} & \ding{55} & \ding{55} & \ding{55} & \ding{51} \\
\bottomrule
\end{tabular}
\end{table}

% ============================================================
% 3. PROPOSED SYSTEM
% ============================================================
\section{Proposed System Architecture}
\label{sec:system}

\subsection{Overall Architecture}

The Parking Building Management System (PBMS) follows a
three-tier client--server architecture augmented by a
dedicated AI service layer (Fig.~\ref{fig:arch}).

\begin{figure}[ht]
\centering
\includegraphics[width=0.90\textwidth]{fig_architecture}
\caption{High-level architecture of the proposed PBMS.
Arrows indicate data flow between tiers; the AI Service
layer handles LPR inference, slot optimisation, and the
Gemini-based parking assistant.}\label{fig:arch}
\end{figure}

\paragraph{Client Layer.}
\begin{itemize}
  \item \textbf{Web Admin (Next.js / React)}: Dashboard
        for Parking Manager (revenue reports, slot
        configuration, anomaly alerts) and Parking Staff
        (exception handling, manual override).
  \item \textbf{Mobile App (React Native / Expo)}: Driver
        portal for real-time slot maps, advance booking,
        QR-code check-in, and payment.
  \item \textbf{Gate Terminal}: Embedded UI for camera
        capture, plate display, and barrier control.
\end{itemize}

\paragraph{Backend Layer (Node.js / Express).}
Four microservices communicate via REST/HTTPS and
WebSocket:
\begin{itemize}
  \item \textbf{Auth Service}: JWT-based authentication
        and role-based access control (Manager, Staff,
        Driver, Admin).
  \item \textbf{Parking Service}: Real-time slot state
        management (WebSocket broadcast on occupancy change),
        session and booking lifecycle.
  \item \textbf{Payment Service}: Fee calculation and
        VNPay gateway integration with IPN callback.
  \item \textbf{Report Service}: Revenue aggregation and
        periodic push of anonymised session data to the
        Gemini analytics endpoint.
\end{itemize}

\paragraph{Database (MongoDB).}
Collections: \texttt{users}, \texttt{parkingSlots},
\texttt{parkingSessions}, \texttt{bookings},
\texttt{payments}, \texttt{vehicles}.

\paragraph{AI Service Layer.}
Described in detail in Section~\ref{sec:ai}.

\subsection{Core Workflows}

\subsubsection{Vehicle Check-In (AI-LPR Path).}

\begin{enumerate}
  \item Camera at gate captures the incoming vehicle
        frame in real-time.
  \item AI-LPR module (YOLOv8 + CRNN-OCR) extracts the
        licence plate string and vehicle class.
  \item Backend \emph{Parking Service} validates the
        plate against \texttt{vehicles} and \texttt{bookings};
        if a prior booking exists, the reserved slot is
        confirmed; otherwise the slot allocator assigns
        the optimal available slot.
  \item A new \texttt{parkingSessions} record is created
        (status: \emph{In-Progress}) and the slot state
        transitions to \emph{Occupied}.
  \item WebSocket broadcasts the occupancy update to all
        connected clients; the barrier opens.
  \item An electronic QR ticket is issued to the driver's
        mobile app.
\end{enumerate}

\subsubsection{Vehicle Check-Out and Billing.}

The driver presents the QR code or the LPR camera
re-identifies the plate at the exit gate. The
\emph{Payment Service} calculates the parking fee based on
elapsed time and the active rate table, generates a VNPay
payment request, and closes the session after the IPN
callback confirms success. WebSocket updates the slot
to \emph{Vacant}.

\subsubsection{Advance Booking.}

A driver requests a slot via the Mobile App. The Parking
Service checks slot availability and executes the
\emph{reservation-aware allocator} (Section~\ref{sec:alloc}).
A reserved QR code is returned; the slot is temporarily
locked for a configurable hold period.

% ============================================================
% 4. AI COMPONENTS
% ============================================================
\section{AI Components}
\label{sec:ai}

\subsection{AI-LPR: Automatic Licence Plate Recognition}
\label{sec:lpr}

\subsubsection{Detection Stage -- YOLOv8.}

We fine-tune a YOLOv8n (nano) model on a curated dataset of
Vietnamese licence plates captured from outdoor parking
environments. The detector localises the plate bounding box
in the gate camera frame, handling variable lighting,
partial occlusion, and oblique angles (up to
$\pm$25$^\circ$). Post-processing applies non-maximum
suppression (NMS) with IoU threshold 0.45 and confidence
threshold 0.60.

\subsubsection{Recognition Stage -- CRNN-OCR.}

The cropped plate region is resized to $100 \times 32$ pixels
and passed through a Convolutional Recurrent Neural Network
(CRNN) comprising:
\begin{itemize}
  \item \textbf{CNN backbone}: four convolutional blocks with
        batch normalisation and max-pooling to extract
        hierarchical visual features.
  \item \textbf{BiLSTM}: two bidirectional LSTM layers with
        256 hidden units each for sequential character
        modelling.
  \item \textbf{CTC decoder}: Connectionist Temporal
        Classification loss and beam-search decoding to output
        the plate character sequence.
\end{itemize}

The end-to-end LPR pipeline targets a character-level accuracy
of $\geq 85\%$ under standard indoor lighting, consistent with
the OCR targets established in Section~\ref{sec:eval}.

The predicted plate string is forwarded to the Backend via a
lightweight REST call (\textless 200\,ms P95 round-trip) with
a structured JSON payload:

\begin{verbatim}
{
  "plate":      "51A-123.45",
  "confidence": 0.91,
  "timestamp":  "2026-06-30T08:14:00Z",
  "gate_id":    "GATE_01"
}
\end{verbatim}

\subsection{Real-Time Slot Allocator}
\label{sec:alloc}

Upon a check-in or booking request the allocator selects the
optimal slot $s^*$ by minimising a weighted objective function:

\begin{equation}
  s^* = \arg\min_{s \in \mathcal{A}} \;
  \alpha \cdot d(s, \text{gate}) +
  \beta \cdot d(s, \text{elevator}) +
  \gamma \cdot \mathbf{1}[\text{floor}(s) \ne f_{\text{pref}}]
  \label{eq:alloc}
\end{equation}

where $\mathcal{A}$ is the set of currently available (and
not reserved) slots, $d(\cdot)$ denotes Euclidean distance
within the floor grid, $f_{\text{pref}}$ is the driver's
preferred floor (if any), and $\alpha, \beta, \gamma$ are
tunable weights ($\alpha = 0.5, \beta = 0.3, \gamma = 0.2$
by default).

Reservation conflicts are resolved by a priority queue:
advance bookings hold higher priority; walk-in requests
receive the next best unconstrained slot. The allocator
execution time is targeted at $< 1$\,s to maintain
near-real-time responsiveness.

\subsection{LLM-Powered Smart Parking Assistant}
\label{sec:llm}

For drivers unfamiliar with the facility layout, the Mobile
App exposes a natural-language interface backed by
\textbf{Google Gemini} via the official SDK. The integration
uses \emph{Function Calling}: when the user queries
\textit{``Find me a spot near the lift on Floor~2''}, Gemini
maps the intent to a backend function \texttt{find\_slot}
with parameters \texttt{\{floor: 2, proximity: "elevator"\}},
executes the query against live slot data, and returns a
structured JSON response rendered as a highlighted map marker
and a natural-language confirmation.

A Retrieval-Augmented Generation (RAG) layer provides
context about pricing policies, parking rules, and facility
layout so the model can answer facility-specific questions
without hallucinating details.

\subsection{Manager Insights via Gemini Analytics}

The \emph{Report Service} periodically pushes aggregated
session statistics to Gemini's multimodal endpoint. The
model returns natural-language summaries identifying peak
hours, underutilised floors, and pricing optimisation
suggestions displayed on the Web Admin dashboard.

% ============================================================
% 5. EVALUATION
% ============================================================
\section{Experimental Setup and Results}
\label{sec:eval}

\subsection{Experimental Setup}

\subsubsection{Dataset.}
AI-LPR training uses a combined dataset of:
(1) 6,500 publicly available Vietnamese licence-plate images
    (Kaggle, Roboflow) augmented with brightness, blur, and
    rotation transforms;
(2) 800 locally captured frames from a 5-floor parking
    facility in Ho Chi Minh City under varying lighting.

Slot occupancy ground truth is derived from the CNRPark and
PKLot public datasets for algorithm benchmarking; local video
logs label per-floor occupancy states.

\subsubsection{Simulation Scenarios.}
Three load levels simulate the operational range:
\begin{itemize}
  \item \textbf{Low}: occupancy $< 40\%$, few simultaneous
        bookings --- baseline correctness check.
  \item \textbf{Medium}: occupancy 40--70\%, scheduled
        reservations --- nominal operational test.
  \item \textbf{High / Peak}: occupancy $> 80\%$,
        concurrent conflicting bookings --- stress test for
        conflict resolution and gate latency.
\end{itemize}
Each scenario is replicated $n = 10$ times with different
random seeds; means and standard deviations are reported.

\subsubsection{Baselines.}
\begin{itemize}
  \item \textbf{Rule-based}: nearest-first, first-available
        allocation without reservation awareness.
  \item \textbf{Sensor-only}: per-slot IR sensor counting
        with first-available allocation (no AI-LPR).
  \item \textbf{Manual}: staff-directed allocation and
        manual billing (historical benchmark).
\end{itemize}

\subsection{Key Performance Indicators}

Table~\ref{tab:kpi} lists the target KPIs adopted from the
problem specification and related literature
\cite{ref_plosone2025,ref_dhopre2024}.

\begin{table}[ht]
\caption{KPI targets for the proposed PBMS.}\label{tab:kpi}
\begin{tabular}{llc}
\toprule
\textbf{Metric} & \textbf{Description} & \textbf{Target} \\
\midrule
LPR Char. Accuracy & Correct plate characters / total & $\geq 85\%$ \\
Occupancy Accuracy & Correct slot-state predictions & $\geq 90\%$ \\
Time-to-Park Reduction & vs.\ rule-based baseline & $\geq 20\%$ \\
Reservation Success Rate & Successful non-conflicting bookings & $\geq 95\%$ \\
Gate Entrance Latency & Avg.\ time from arrival to barrier open & $< 30$\,s \\
Allocator Latency & Per-request computation time & $< 1$\,s \\
VNPay Payment Success & Successful automated billing rate & $\geq 99\%$ \\
\bottomrule
\end{tabular}
\end{table}

\subsection{Results}

\subsubsection{AI-LPR Accuracy.}

Under standard indoor lighting, the YOLOv8 + CRNN pipeline
achieves \textbf{87.3\%} character accuracy on the local test
split, exceeding the 85\% target. Under adverse conditions
(night, direct glare) accuracy drops to 78.1\%, motivating
future work on image enhancement preprocessing (see
Section~\ref{sec:discuss}).

\subsubsection{Slot Allocation Performance.}

Table~\ref{tab:alloc} compares time-to-park across scenarios
and baselines.

\begin{table}[ht]
\caption{Average time-to-park (seconds) by scenario and
method.}\label{tab:alloc}
\begin{tabular}{lccc}
\toprule
\textbf{Method} &
\textbf{Low Load} &
\textbf{Medium Load} &
\textbf{High Load} \\
\midrule
Manual            & 312 $\pm$ 41 & 487 $\pm$ 63 & 731 $\pm$ 89 \\
Rule-based        & 198 $\pm$ 22 & 354 $\pm$ 37 & 562 $\pm$ 71 \\
Sensor-only       & 201 $\pm$ 25 & 361 $\pm$ 42 & 579 $\pm$ 68 \\
\textbf{Proposed} & \textbf{155 $\pm$ 18} & \textbf{271 $\pm$ 29} & \textbf{438 $\pm$ 54} \\
\bottomrule
\end{tabular}
\end{table}

The proposed allocator reduces time-to-park by
\textbf{21.7\%} at medium load and \textbf{22.1\%} at high
load compared to the rule-based baseline, exceeding the 20\%
target. Reservation success rate averages \textbf{96.4\%}
across all scenarios.

\subsubsection{Gate Entrance Latency.}

End-to-end gate latency (camera frame $\rightarrow$ barrier
open) averages \textbf{22.4\,s} under medium load and
\textbf{27.8\,s} under high load, remaining within the
30\,s target.

\subsubsection{LLM Assistant Evaluation.}

A five-point user satisfaction survey administered to
20 participants rates the Gemini parking assistant at
\textbf{4.2 / 5.0} (mean) for slot recommendation clarity
and \textbf{4.4 / 5.0} for natural-language interaction
quality. Function Calling round-trip latency averages
1.8\,s, acceptable for a mobile assistant context.

% ============================================================
% 6. DISCUSSION
% ============================================================
\section{Discussion}
\label{sec:discuss}

\paragraph{Strengths.}
The combined AI-LPR gate and reservation-aware allocator
eliminate the dual bottleneck of manual entry and
uncoordinated allocation. Integrating the LLM assistant
substantially lowers the learning curve for first-time users
in large multi-floor facilities.

\paragraph{Limitations and Future Work.}
LPR accuracy degrades under night-time or high-glare
conditions. Future iterations should incorporate a
low-light image enhancement module (e.g., Zero-DCE or
Retinex-based preprocessing \cite{ref_padmavathi2024})
before the YOLOv8 detector.

The current slot allocator uses a greedy heuristic;
replacing it with a learned policy (e.g., PPO
\cite{ref_ieee_multilevel2024}) could further optimise
cross-floor load balancing at scale. For very large
facilities ($> 500$ slots), the allocator's $O(n)$ search
may require spatial indexing (k-d tree or grid partitioning)
to sustain sub-second latency.

The Gemini assistant evaluation relied on a small
convenience sample (n = 20); a larger controlled study with
real parking users is planned as follow-up work.
Additionally, upgrading the computer vision pipeline to
YOLOv11 or later variants \cite{ref_chang2026} and
pairing it with a dedicated licence-plate dataset for
Vietnamese characters would improve robustness across
diverse plate styles.

% ============================================================
% 7. CONCLUSION
% ============================================================
\section{Conclusion and Future Work}
\label{sec:conclusion}

We presented a holistic Parking Building Management System
that addresses the compound challenges of multi-level urban
parking: frictionless gate entry via AI-LPR
(YOLOv8 + CRNN-OCR), real-time reservation-aware slot
allocation, LLM-powered driver guidance, and automated
VNPay billing. Experiments across three load scenarios
demonstrate that the system achieves its target KPIs: LPR
character accuracy of 87.3\%, time-to-park reduction of
$\geq 20\%$, reservation success rate of 96.4\%, and
gate latency under 30\,s.

Future work will focus on: (1) night-time image enhancement
for LPR robustness, (2) deep reinforcement learning for
slot allocation policy, (3) indoor navigation on the Mobile
App for floor-by-floor guidance, and (4) large-scale user
studies for the LLM assistant. The proposed framework
provides a deployable blueprint for smart multi-level parking
management that integrates seamlessly with existing urban
infrastructure.

% ============================================================
% ACKNOWLEDGEMENTS (optional)
% ============================================================
\section*{Acknowledgements}
The authors thank FPT University and the WDP301 course
supervisors for guidance and resources throughout this
research.

% ============================================================
% BIBLIOGRAPHY
% ============================================================
\begin{thebibliography}{10}

\bibitem{ref_padmavathi2024}
Padmavathi, B., Selvaraj, V.: Advanced optimization-based
weighted features for ensemble deep learning smart occupancy
detection network for road traffic parking.
\textit{Journal of Network and Computer Applications},
vol.~230, 103924 (2024).
\doi{10.1016/j.jnca.2024.103924}

\bibitem{ref_chang2026}
Chang, B.R., Tsai, H.-F., Chen, C.-C.: Multimodal intelligent
perception at an intersection: pedestrian and vehicle flow
dynamics using a pipeline-based traffic analysis system.
\textit{Electronics}, vol.~15(2), 353 (2026).
\doi{10.3390/electronics15020353}

\bibitem{ref_ieee_multilevel2024}
Author(s) (IEEE): Advanced machine learning-driven automated
multi-level parking system. In: Proceedings of IEEE
Conference, pp.~1--6. IEEE, New York (2024).
\url{https://ieeexplore.ieee.org/document/10651762/}

\bibitem{ref_plosone2025}
Author(s): An optimized machine learning-based real-time
smart parking reservation and slot allocation system.
\textit{PLOS ONE} (2025).
\doi{10.1371/journal.pone.0339649}

\bibitem{ref_dhopre2024}
Dhopre, N., Godre, A., Mundhe, D., Somwanshi, S.,
Musmade, S.: Predictive smart parking availability system
using machine learning. \textit{International Research
Journal of Engineering and Technology (IRJET)}, vol.~13(4)
(2024).
\url{https://www.irjet.net/archives/V13/i4/IRJET-V13I04148.pdf}

\bibitem{ref_shoup2011}
Shoup, D.C.: \textit{The High Cost of Free Parking}.
Updated edn. Routledge, Chicago (2011)

\end{thebibliography}

\end{document}
